\documentclass[11pt,a4paper]{article}
\usepackage[margin=0.82in,headheight=14pt,footskip=26pt]{geometry}
\usepackage[T1]{fontenc}
\usepackage[utf8]{inputenc}
\usepackage{lmodern}
\usepackage{microtype}
\usepackage{amsmath,amssymb,amsthm,mathtools}
\usepackage{graphicx}
\usepackage{booktabs,tabularx,array}
\usepackage{natbib}
\usepackage{xcolor}
\usepackage{hyperref}
\usepackage[nameinlink,noabbrev]{cleveref}
\usepackage[section]{placeins}
\usepackage{fancyhdr}
\usepackage{enumitem}
\definecolor{Ink}{HTML}{172B3A}
\definecolor{Accent}{HTML}{007C91}
\definecolor{Muted}{HTML}{526574}
\hypersetup{colorlinks=true,linkcolor=Accent,citecolor=Accent,urlcolor=Accent,pdftitle={Pulse-Resolved Identifiability: An Adaptive Design Principle for Separating Drift and Diffusion in Destructive Single-Cell Assays},pdfauthor={Devanik Debnath}}
\pagestyle{fancy}
\fancyhf{}
\fancyhead[L]{\small\textcolor{Muted}{Pulse-resolved identifiability}}
\fancyhead[R]{\small\textcolor{Muted}{Research preprint · 24 October 2026}}
\fancyfoot[C]{\small\thepage}
\renewcommand{\headrulewidth}{0.35pt}
\renewcommand{\footrulewidth}{0pt}
\usepackage{titlesec}
\titleformat{\section}{\Large\bfseries\color{Ink}}{\thesection}{0.65em}{}
\titleformat{\subsection}{\large\bfseries\color{Ink}}{\thesubsection}{0.65em}{}
\titleformat{\subsubsection}{\normalsize\bfseries\color{Ink}}{\thesubsubsection}{0.65em}{}
\setlength{\parindent}{1em}
\setlength{\parskip}{0.22em}
\setlist{nosep,leftmargin=1.5em}
\newtheorem{proposition}{Proposition}
\newtheorem{corollary}{Corollary}
\newtheorem{definition}{Definition}
\newcommand{\E}{\mathbb{E}}
\newcommand{\Var}{\mathrm{Var}}
\newcommand{\R}{\mathbb{R}}
\newcommand{\FIM}{\mathcal{I}}
\newcommand{\diag}{\mathrm{diag}}
\newcommand{\logdet}{\mathrm{logdet}}
\title{\vspace{-1.4em}\textbf{\LARGE Pulse-Resolved Identifiability}\\[0.35em]\large An Adaptive Design Principle for Separating Drift and Diffusion in Destructive Single-Cell Assays}
\author{\textbf{Devanik Debnath (Devanik21)}\\
\small Department of Electronics and Communication Engineering\\
\small National Institute of Technology Agartala, India\\
\small ECE '26 \quad\textbullet\quad ORCID: \href{https://orcid.org/0009-0001-8310-6921}{0009-0001-8310-6921}}
\date{\small Research preprint \quad|\quad 24 October 2026}
\begin{document}
\maketitle
\vspace{-0.8em}
\begin{abstract}
Single-cell assays often destroy the measured cell, so temporal experiments yield unpaired population snapshots rather than trajectories of the same individuals. Distribution-matching and stochastic-trajectory methods can extract useful structure from such data, but a fitted trajectory is not automatically identifiable from the observations. We isolate a basic ambiguity in a controlled Ornstein--Uhlenbeck model: a stationary variance identifies only a drift--diffusion ratio, and one post-pulse population mean supplies too few independent moment constraints to identify relaxation, input gain, and noise intensity simultaneously. We derive an exact pulse-response relation showing that two distinct post-pulse mean measurements identify the relaxation rate under explicit model assumptions; a calibrated observation-noise variance then permits recovery of input gain and diffusion intensity. We turn this relation into an information-aware sampling rule that greedily selects measurement times using a local Fisher-information objective while respecting a state-amplitude constraint. A reproducible synthetic benchmark compares this rule with uniform, early-only, and random schedules, reports parameter error and information geometry, and includes tests designed to expose confounding rather than to guarantee a preferred outcome. The contribution is deliberately bounded: this is a worked identifiability certificate and design benchmark, not a new general theory of optimal experimental design, a validated biological model, or evidence of rejuvenation or clinical efficacy. The framework is intended to help researchers decide when cross-sectional single-cell measurements can support kinetic claims and when additional calibration or intervention is required.
\end{abstract}
\noindent\textbf{Keywords:} single-cell perturbation; structural identifiability; stochastic differential equations; Fisher information; active experimental design; destructive sampling; optimal design; scientific machine learning
\section{Research question and contribution}
A central practical tension in single-cell biology is that high-dimensional molecular measurements are frequently destructive, while many scientific questions concern dynamics: how rapidly a state relaxes, how strongly an intervention shifts it, and how much variation is generated by stochastic processes. Cross-sectional snapshots can be collected at multiple times, but the cells at time $t_1$ are not the same individuals measured at $t_2$. Methods such as Waddington-OT, CellOT, and newer stochastic transport models use distributional structure to infer trajectories or perturbation responses \citep{schiebinger2019,bunne2023,liu2026}. Those methods do not remove the need to ask a prior question: which parameters are determined by the available observations, and which remain confounded?
This paper studies that question in a deliberately small model whose assumptions can be audited. The goal is not to propose a universal cell-state equation. It is to provide a transparent test case in which the distinction between fitting a distribution and identifying its generating kinetics can be derived exactly, tested numerically, and translated into an experimental-design rule.
The contributions are fourfold:
\begin{enumerate}
\item We derive the moment-level confounding of relaxation rate, input gain, diffusion intensity, and measurement noise in a stationary linear stochastic model.
\item We give an exact post-pulse ratio identity and state the conditions under which it identifies the relaxation rate from unpaired population means.
\item We implement a sequential, information-aware sampling policy that chooses time points from a candidate grid using local Fisher information and a declared amplitude constraint.
\item We release a synthetic benchmark with deterministic seeds, baseline schedules, failure tests, generated plots, and machine-readable outputs. No biological dataset is claimed to have been processed.
\end{enumerate}
The novelty claim is intentionally narrow. Optimal design, Fisher-information analysis of stochastic kinetics, and single-cell perturbation modelling are established areas \citep{komorowski2011,fox2020,iolov2017,bunne2023}. We do not claim that D-optimal design or pulse-based system identification is new. The contribution is an executable, explicitly scoped bridge between these literatures for the destructive-snapshot setting, including a precise statement of the calibration assumptions required for separating drift and diffusion.
\section{Literature position and epistemic scope}
Optimal-transport approaches have inferred developmental trajectories from time-stamped single-cell populations, while CellOT learns perturbation-conditioned maps between unpaired distributions \citep{schiebinger2019,bunne2023}. The scPerturb resource was assembled to improve interoperability and benchmarking across perturbation-response datasets \citep{peidli2024}; pseudotime work has also explicitly treated confounding in cross-sectional single-cell time series \citep{reid2016}. DiffusionOT extends snapshot-based inference by modelling stochastic effects in a continuous-time framework \citep{liu2026}; the October 2026 CloneOT preprint uses partial lineage information to learn perturbation-specific transport costs \citep{chen2026}. These works establish that distributional modelling is active and rapidly evolving. They do not imply that all latent parameters are identifiable under every sampling design.
The complementary design literature is also substantial. Fisher-information methods have long been used to analyse sensitivity and identifiability in stochastic chemical kinetics \citep{komorowski2011}. Fox, Neuert, and Munsky demonstrated model-driven allocation of destructive single-cell measurements across time in a fluctuating environment \citep{fox2020}. Iolov, Ditlevsen, and Longtin studied input design for an Ornstein--Uhlenbeck process observed through threshold-crossing times \citep{iolov2017}. Our formulation therefore sits within existing system-identification and optimal-design traditions. The comparison in \cref{tab:literature} makes the boundary explicit: the present work does not claim to invent information-based sampling, but isolates a pulse-response identity for destructive population snapshots.
\begin{table}[!htbp]
\centering
\small
\caption{Relevant literature and the specific boundary of the present contribution.}
\label{tab:literature}
\begin{tabularx}{\linewidth}{@{}p{2.35cm}p{2.25cm}X X@{}}
\toprule
Work & Data or model setting & Established capability & Boundary relevant here\\
\midrule
Schiebinger et al. (2019) & Time-stamped single-cell RNA-seq & Distributional transport for developmental trajectories & Does not make every kinetic parameter identifiable from any sampling schedule\\
Bunne et al. (2023) & Unpaired control/perturbed cells & Neural optimal-transport maps for perturbation response & Individual-cell maps are modelling objects when observations are unpaired\\
Peidli et al. (2024) & Perturbation-response resource & Harmonized datasets and distributional benchmarks & Dataset harmonization alone does not resolve structural confounding\\
Liu et al. (2026) & Multi-time-point snapshots & Stochastic dynamics and cell-fate landscape inference & Inference quality remains conditional on model and sampling assumptions\\
Fox et al. (2020) & Stochastic gene regulation & Information-based allocation of single-cell samples over time & Provides a close precedent for design; this paper emphasizes an exact drift/diffusion identifiability certificate\\
Iolov et al. (2017) & Controlled OU process, threshold times & Information-based perturbation design & Different observation model and objective; the present model uses population means and variances\\
\bottomrule
\end{tabularx}
\end{table}
\subsection{P1--P5 classification}
The exact solution of the stated linear stochastic differential equation and the algebraic confounding statements are P1 mathematical consequences. The Fisher-information objective and the greedy schedule are P2 design choices, not a theorem of biological optimality. Any future claim that the same schedule improves prediction across cell types or tissues would be P3 and would require external validation. The mapping from this scalar model to a high-dimensional cellular state is P4 structural analogy only. Metaphors about cells “remembering” pulses or “escaping” a landscape are P5 and are not used as evidence.
\section{Model, observables, and assumptions}
Let $x(t)\in\R$ denote a single latent state coordinate for a cell or cell-like unit. Under a controlled pulse, assume
\begin{equation}
dx(t)=\big[-\kappa x(t)+\beta u(t)\big]dt+\sqrt{q}\,dW_t,
\label{eq:sde}
\end{equation}
where $\kappa>0$ is the relaxation rate, $\beta>0$ is the input gain, $q>0$ is the diffusion intensity, and $W_t$ is standard Brownian motion. The applied input is a rectangular pulse
\begin{equation}
u(t)=
\begin{cases}
u_0,&0\le t\le \tau,\\
0,&t>\tau,
\end{cases}
\label{eq:pulse}
\end{equation}
with known amplitude $u_0\ne0$ and duration $\tau>0$. We assume the initial state is drawn from the stationary distribution of the uncontrolled process, $x(0)\sim\mathcal{N}(0,q/(2\kappa))$. At each scheduled time $t_j$, a different randomized cohort is measured:
\begin{equation}
y_{ij}=x_i(t_j)+\epsilon_{ij},\qquad
\epsilon_{ij}\sim\mathcal{N}(0,r),
\label{eq:observation}
\end{equation}
where $r\ge0$ is the measurement-noise variance. The errors and Brownian paths are independent across cells; the pulse, initial distribution, and observation protocol are assumed comparable across cohorts. The principal analysis treats $r$ as known from an independent calibration. We later show why this assumption matters. The notation and scope of the assumptions are summarized in \cref{tab:notation}.
\begin{table}[!htbp]
\centering
\small
\caption{Notation, observables, and model assumptions.}
\label{tab:notation}
\begin{tabularx}{\linewidth}{@{}p{1.65cm}p{2.4cm}X@{}}
\toprule
Symbol & Meaning & Assumption or observable status\\
\midrule
$\kappa$ & Relaxation rate & Unknown positive parameter; time unit is inverse to $\kappa$\\
$\beta$ & Pulse-to-state gain & Unknown positive parameter; units depend on state and input scaling\\
$q$ & Diffusion intensity & Unknown positive parameter; state-squared per time\\
$u_0,\tau$ & Pulse amplitude and duration & Set by experiment; not inferred in the base model\\
$r$ & Observation-noise variance & Known in the main identifiability result; unknown $r$ creates extra confounding\\
$m(t)$ & Population mean after pulse & Estimated from destructive cohorts; cells are not paired across times\\
$v$ & Observed stationary variance & $q/(2\kappa)+r$ under stationary initialisation\\
$n_j$ & Cells sampled at time $t_j$ & Independent biological units is an idealization; donor and batch structure are omitted\\
\bottomrule
\end{tabularx}
\end{table}
\subsection{Closed-form moments}
The mean $m(t)=\E[x(t)]$ satisfies $\dot m=-\kappa m+\beta u(t)$ with $m(0)=0$. Therefore
\begin{equation}
m(t)=
\begin{cases}
\dfrac{\beta u_0}{\kappa}\big(1-e^{-\kappa t}\big),&0\le t\le\tau,\\[0.65em]
\dfrac{\beta u_0}{\kappa}\big(1-e^{-\kappa\tau}\big)e^{-\kappa(t-\tau)},&t>\tau.
\end{cases}
\label{eq:mean}
\end{equation}
Because the input is deterministic and affects only the mean, the centered process remains an Ornstein--Uhlenbeck process. Stationary initialisation makes its latent variance constant, so the observed variance is
\begin{equation}
v=\Var[y(t)]=\frac{q}{2\kappa}+r.
\label{eq:variance}
\end{equation}
Equation \cref{eq:variance} is the core ambiguity: without a separate estimate of $\kappa$, variance alone identifies only a combination of $q$ and $\kappa$; if $r$ is also unknown, even that combination is confounded with measurement noise.
\section{Identifiability results}
\begin{definition}[Moment identifiability]
A parameter vector is locally moment-identifiable under a design if the mapping from that vector to the collection of observed population means and variances has full column rank in its Jacobian at the parameter point of interest. This is a local statement and does not by itself guarantee stable finite-sample estimation.
\end{definition}
\begin{proposition}[Stationary variance confounding]
If only the stationary observed variance $v$ is available and $r$ is known, then infinitely many positive pairs $(\kappa,q)$ produce the same variance. If $r$ is unknown, the confounding set is larger.
\end{proposition}
\begin{proof}
For any chosen $\kappa>0$, set $q=2\kappa(v-r)$ whenever $v>r$. Then \cref{eq:variance} holds. Thus the observation determines $q/\kappa=2(v-r)$, not the two parameters separately. If $r$ is unknown, any admissible decomposition of $v$ into $q/(2\kappa)+r$ is observationally equivalent at the variance level.
\end{proof}
\begin{proposition}[Single-time underdetermination]
With known $r$, one post-pulse population mean and one stationary variance cannot identify all three parameters $(\kappa,\beta,q)$ from moments alone.
\end{proposition}
\begin{proof}
The single-time mean contributes one scalar equation involving $\kappa$ and $\beta$; the stationary variance contributes one scalar equation involving $\kappa$ and $q$. The moment map has two outputs for three unknowns, so its Jacobian has rank at most two. In particular, varying $\kappa$ can be compensated by changing $\beta$ to preserve the mean and changing $q$ to preserve the variance, locally wherever the derivatives are nonzero.
\end{proof}
\begin{proposition}[Two post-pulse means identify the relaxation rate]
Let $t_1,t_2>\tau$ with $t_1\ne t_2$, and assume the mean response is nonzero and measured without systematic cohort bias. Under \cref{eq:sde}--\cref{eq:pulse}, the two post-pulse population means satisfy
\begin{equation}
\kappa=\frac{\log m(t_1)-\log m(t_2)}{t_2-t_1},
\label{eq:kappa_ratio}
\end{equation}
for $m(t_1),m(t_2)>0$ (or for consistently signed nonzero means after taking the log of their magnitude).
\end{proposition}
\begin{proof}
For $t>\tau$, \cref{eq:mean} has the form $m(t)=C e^{-\kappa(t-\tau)}$, where $C=(\beta u_0/\kappa)(1-e^{-\kappa\tau})$ is constant across observation times. Hence $m(t_1)/m(t_2)=e^{\kappa(t_2-t_1)}$. Taking logarithms gives \cref{eq:kappa_ratio}.
\end{proof}
\begin{corollary}[Recovering gain and diffusion under calibration]
If $\kappa$ is identified from two distinct post-pulse means, $u_0$ and $\tau$ are known, and $r$ is calibrated, then $\beta$ and $q$ are identified from
\begin{equation}
\beta=\frac{\kappa\,m(t)e^{\kappa(t-\tau)}}{u_0(1-e^{-\kappa\tau})},\qquad
q=2\kappa(v-r).
\label{eq:recover}
\end{equation}
\end{corollary}
These identities concern the idealized moment map. They do not guarantee practical identifiability when the response has decayed below the measurement floor, when the pulse is weak, when cohorts differ systematically, or when the model is misspecified. The observation regimes and their limits are compared in \cref{tab:identifiability}. The exact result suggests that useful schedules should include separated post-pulse times but should avoid choosing only very late times with negligible signal.
\begin{table}[!htbp]
\centering
\small
\caption{Identifiability under different observation designs in the stated model.}
\label{tab:identifiability}
\begin{tabularx}{\linewidth}{@{}p{2.4cm}p{2.2cm}X p{2.0cm}@{}}
\toprule
Design & Independent moment constraints & What can be identified & Main failure mode\\
\midrule
Unperturbed stationary snapshots & Mean $0$; variance $v$ & $q/\kappa$ if $r$ is known & Drift--diffusion ratio confounding\\
One post-pulse time & One mean and one variance & At most two independent combinations of $(\kappa,\beta,q)$ & Under-determined three-parameter map\\
Two distinct post-pulse times + variance & Two means and one variance & $\kappa$, then $\beta$, then $q$ if $r$ known & Low signal or systematic cohort shift\\
Same design, unknown $r$ & Two means and observed variance & $\kappa$ and $\beta$; $q$ remains confounded with $r$ & Noise calibration missing\\
Many time points, wrong dynamics & More moment samples, misspecified model & Parameters of the fitted surrogate only & Good fit mistaken for mechanism\\
\bottomrule
\end{tabularx}
\end{table}
\section{Fisher information and adaptive sampling}
For $\theta=(\log\kappa,\log\beta,\log q)$, the observed state at a scheduled time is Gaussian with mean $m(t;\theta)$ and variance $v(\theta)$. For $n$ independent cells sampled at time $t$, the expected Fisher information is
\begin{equation}
\FIM_t(\theta)=n\left[
\frac{\nabla_\theta m(t)\nabla_\theta m(t)^\top}{v}
+\frac{1}{2}\nabla_\theta\log v\,\nabla_\theta\log v^\top
\right].
\label{eq:fim}
\end{equation}
For a schedule $\mathcal{T}=\{t_1,\ldots,t_K\}$ with $n$ cells at each time, information adds across independent cohorts:
\begin{equation}
\FIM_{\mathcal{T}}(\theta)=\sum_{j=1}^{K}\FIM_{t_j}(\theta).
\label{eq:fim_sum}
\end{equation}
The design policy chooses a new time from a finite candidate set $\mathcal{G}$ to increase a regularized log-determinant objective:
\begin{equation}
t_{k+1}=\arg\max_{t\in\mathcal{G}\setminus\mathcal{T}_k}
\log\det\left[\FIM_{\mathcal{T}_k}(\widehat\theta_k)+\FIM_t(\widehat\theta_k)+\lambda I\right],
\label{eq:greedy}
\end{equation}
subject to a declared amplitude constraint $\max_{t\in\mathcal{G}}|m(t;\widehat\theta_k)|\le M_{\max}$ and a fixed sampling budget. The estimate $\widehat\theta_k$ is updated from the population moments collected so far using a regularized Gaussian likelihood. This is a local greedy design, not a globally optimal policy and not a learned biological law. The small ridge $\lambda I$ makes the objective finite before the information matrix becomes full rank; it must not be interpreted as observed information.
\subsection{Algorithmic interpretation}
The adaptive component is an active-learning policy: at each round it selects the measurement time expected to reduce local parameter uncertainty, observes a new independent cohort, updates its parameter estimate, and recomputes the information geometry. The benchmark conditions and sampling budget are specified in \cref{tab:protocol}. The policy does not need to match individual cells across time. It does require randomized or otherwise comparable cohorts and a credible observation model. If donor effects, cell-state composition, or batch drift violate these assumptions, the apparent information gain can be spurious.
\begin{table}[!htbp]
\centering
\small
\caption{Synthetic benchmark protocol and comparison conditions.}
\label{tab:protocol}
\begin{tabularx}{\linewidth}{@{}p{2.4cm}X@{}}
\toprule
Item & Prespecified setting\\
\midrule
Latent model & Scalar controlled Ornstein--Uhlenbeck process, stationary initial distribution\\
Pulse & Rectangular input with fixed amplitude and duration; pulse values are normalized model units\\
Unknown parameters & Log relaxation rate, log input gain, and log diffusion intensity\\
Observation & Independent Gaussian cross-sectional samples with known additive measurement-noise variance\\
Candidate times & Fixed grid spanning the pulse-on, early post-pulse, and late post-pulse periods\\
Policies & Sequential greedy D-optimal, uniform, early-only, and random schedules\\
Evaluation & Monte Carlo parameter estimation under a declared log-parameter prior around a nominal model\\
Randomness & Fixed pseudorandom seeds; no external dataset or network request\\
Interpretation & Synthetic method stress test only; no empirical biological or clinical validation\\
\bottomrule
\end{tabularx}
\end{table}
\subsection{Why the AI component is bounded}
A neural model could learn a richer posterior or policy from previous experiments, but such a model would introduce additional assumptions and would be difficult to validate without real assay data. The present release uses a transparent greedy policy because its decision is inspectable and its information objective is independently recomputable. The useful AI capability is sequential experimental planning under uncertainty, not autonomous discovery of an unobserved biological mechanism. A future learned policy should be compared against this simple baseline and should be rejected if it fails to improve held-out parameter recovery or violates the amplitude constraint.
\section{Numerical benchmark}
All numerical results in this release are generated by \texttt{verify\_model.py}. The data are synthetic, drawn from the stated Gaussian observation model, and are archived under \texttt{results/}. No public single-cell dataset was downloaded or analyzed. The fixed seed makes the current benchmark reproducible; it does not establish robustness to arbitrary seeds, model classes, or biological heterogeneity.
The controlled mean response in \cref{fig:mean} shows how the post-pulse slope depends on relaxation. The variance-only likelihood ridge in \cref{fig:ridge} exposes the drift--diffusion ambiguity. The schedule comparison in \cref{fig:information} reports the local information objective, while \cref{fig:rmse} reports finite-sample log-parameter error and bootstrap intervals. The dependency graph in \cref{fig:dependency} is derived from the equations, the coupling graph in \cref{fig:coupling} is derived from the Fisher-information matrix, and the two flow diagrams in \cref{fig:pipeline,fig:loop} document the assay and adaptive-design logic. The output table is generated from the current execution rather than copied from a prior run; the actual values should be read together with the seed, sample size, and estimator described in the results record.
\begin{figure}[!htbp]
\centering
\includegraphics[width=0.93\linewidth]{figures/fig1_mean_response.pdf}
\caption{Pulse-response curves for the synthetic linear model. Points and error bars are generated from simulated independent cohorts; curves show the corresponding model means. The plot is an illustration of identifiability structure, not measured cell data.}
\label{fig:mean}
\end{figure}
\begin{figure}[!htbp]
\centering
\includegraphics[width=0.93\linewidth]{figures/fig2_variance_confounding.pdf}
\caption{Variance-only likelihood surface for synthetic stationary observations. The ridge follows near-constant $q/(2\kappa)$, illustrating why stationary variance alone cannot separate drift from diffusion.}
\label{fig:ridge}
\end{figure}
\begin{figure}[!htbp]
\centering
\includegraphics[width=0.93\linewidth]{figures/fig3_information_schedule.pdf}
\caption{Regularized log-determinant Fisher information for candidate schedules at the nominal parameter vector. The random bar summarizes a deterministic set of randomly drawn schedules; it is not a biological benchmark.}
\label{fig:information}
\end{figure}
\begin{figure}[!htbp]
\centering
\includegraphics[width=0.93\linewidth]{figures/fig4_parameter_rmse.pdf}
\caption{Monte Carlo root-mean-square error in log parameters for the adaptive, uniform, early-only, and random policies. Error bars are bootstrap intervals over the recorded synthetic replicates. The ordering is an empirical result of this benchmark, not an asserted theorem.}
\label{fig:rmse}
\end{figure}
\begin{figure}[!htbp]
\centering
\includegraphics[width=0.92\linewidth]{figures/fig5_model_dependency_graph.pdf}
\caption{Structural dependency graph of the equations. Edges indicate mathematical dependence in the assumed model; they are not inferred gene-regulatory or causal edges in a biological system.}
\label{fig:dependency}
\end{figure}
\begin{figure}[!htbp]
\centering
\includegraphics[width=0.82\linewidth]{figures/fig6_information_coupling_graph.pdf}
\caption{Information-coupling graph derived from the normalized Fisher-information matrix for the nominal greedy schedule. Edge widths represent absolute local parameter coupling; they do not imply causal interaction between biological parameters.}
\label{fig:coupling}
\end{figure}
\begin{figure}[!htbp]
\centering
\includegraphics[width=0.98\linewidth]{figures/fig7_assay_pipeline.pdf}
\caption{Analysis pipeline from pulse specification to identifiability assessment. Each stage produces an explicit object that can be inspected or falsified.}
\label{fig:pipeline}
\end{figure}
\begin{figure}[!htbp]
\centering
\includegraphics[width=0.98\linewidth]{figures/fig8_adaptive_loop.pdf}
\caption{Sequential adaptive-design loop. The policy updates its current parameter estimate and chooses a new time under a declared sampling budget and amplitude constraint.}
\label{fig:loop}
\end{figure}
\FloatBarrier
\subsection{Numerical results and interpretation}
\input{results/generated_results.tex}
The current-run numerical summary is reported in \cref{tab:results}. The benchmark is designed to expose trade-offs, not to guarantee that the adaptive policy wins. The greedy policy can perform poorly when the nominal parameter vector is far from the true system, when the mean response is near zero, or when the Gaussian model is false. Early-only schedules may estimate the gain but have little leverage on the post-pulse decay; late-only schedules may span time widely but have weak signal. The most informative design balances parameter sensitivity with measurement precision, which is why a log-determinant score and a state-amplitude constraint are reported separately.
\section{Falsification, robustness, and competing explanations}
The central claims can be challenged in at least five ways. First, measurement-noise calibration may be inaccurate; then the inferred diffusion intensity can absorb observation noise. Second, cohort differences can create a time-dependent mean shift that resembles exponential relaxation. Third, the true dynamics may be nonlinear, multi-stable, or state-dependent in their diffusion. Fourth, cell proliferation, death, or state composition can change population moments without any within-cell relaxation. Fifth, a pulse may alter variance or the initial distribution, invalidating the stationary-variance simplification.
\begin{table}[!htbp]
\centering
\small
\caption{Failure modes and observations that would falsify or narrow the proposed design principle.}
\label{tab:falsification}
\begin{tabularx}{\linewidth}{@{}p{2.2cm}X X@{}}
\toprule
Failure mode & Diagnostic or negative control & Consequence for interpretation\\
\midrule
Unknown measurement noise & Re-estimate $r$ from calibration controls; vary $r$ across plausible bounds & Do not claim $q$ is identified if its estimate tracks $r$\\
Cohort or batch shift & Randomize time-point batches; include batch covariates; replicate donors & A time trend that disappears after adjustment is not kinetic evidence\\
Weak pulse response & Inspect signal-to-noise and confidence intervals for post-pulse means & The exact ratio may be numerically unstable even when structurally valid\\
Non-exponential decay & Compare residuals against multi-exponential and nonlinear alternatives & Report surrogate parameters only; do not interpret $\kappa$ as a unique mechanism\\
Heterogeneous cell populations & Stratify by cell type and donor; check mixture models and composition changes & Population means may mix within-state dynamics with changing proportions\\
Design-policy overfit & Evaluate on held-out parameter draws, random seeds, and alternate noise levels & Reject the policy if its apparent gain does not generalize\\
\bottomrule
\end{tabularx}
\end{table}
The concrete failure modes, diagnostics, and falsification consequences are summarized in \cref{tab:falsification}. A practical experiment should include negative controls, independent biological replicates, and predeclared model checks. If only one post-pulse time is affordable, the honest output is a partially identified parameter set, not a single point estimate with unjustified precision. If the noise calibration is unavailable, the appropriate result is an identified combination of parameters or a sensitivity range over $r$. If an alternative model explains the same observations equally well, the data do not establish a unique biological mechanism.
\section{Human benefit, safety, and translation boundary}
Better experimental design can reduce unnecessary sampling and help researchers distinguish intervention effects from stochastic variation. In regenerative biology, that could improve the prioritization of follow-up assays for cell-state control, repair, or reprogramming. The causal path to human benefit is indirect: a design rule might improve parameter estimation, which might improve the choice of later experiments, which could eventually inform a validated biological model. Each step requires evidence. This paper provides no intervention dose, treatment recommendation, clinical protocol, or claim of therapeutic efficacy.
The model omits donor-level random effects, cell-cycle state, cell death, nonlinear feedback, transcriptional bursting, latent cell-type mixtures, and non-Gaussian measurement errors. These omissions are material. The scalar state $x$ is not a molecular marker, biological age, epigenetic clock, or clinically meaningful outcome. The pulse amplitude bound is a normalized mathematical constraint, not a safe biological dose. No human or animal experiment was performed or proposed as part of this computational release.
\section{Reproducibility and limitations}
The release includes the exact LaTeX source, bibliography, executable verification script, generated figure assets, source data for the plots, environment report, and a checksum manifest. Reproduce the synthetic benchmark from the release directory with:
\begin{verbatim}
python -m py_compile verify_model.py deploy.py
python verify_model.py
python deploy.py --self-test
python deploy.py --validate
\end{verbatim}
The build uses Python with NumPy, SciPy, Matplotlib, NetworkX, pypdf, and PyYAML, plus a LaTeX toolchain with BibTeX and Poppler utilities for PDF checks. Exact versions are recorded in \texttt{results/environment_report.json}. The numerical tests verify the equations and outputs of this implementation; they do not validate the biological relevance of the scalar model. The generated figures use synthetic model data and the graph diagrams are conceptual or equation-derived. There is no external dataset, wet-lab evidence, clinical validation, or DOI deposit associated with this release.
\section{Conclusion}
A population-level fit is not the same as parameter identifiability. In the controlled linear model, stationary variance confounds diffusion with relaxation; one post-pulse mean leaves a three-parameter moment map underdetermined; and two distinct post-pulse means yield an exact relaxation-rate identity when the signal is nonzero and the pulse protocol is known. With calibrated measurement noise, the remaining gain and diffusion parameters follow from closed-form relations. A local Fisher-information planner can use these facts to allocate destructive measurements across time, but its performance is model- and noise-dependent. The next meaningful test is a pre-registered evaluation on a public perturbation time course with donor-aware baselines and explicit comparison to existing methods. Until that test is completed, this release should be read as a reproducible mathematical benchmark and design hypothesis, not as empirical evidence about cell fate.
\section*{Research-process credit and AI disclosure}
Co-equal research-process credit is given to Devanik Debnath and Xylia, the autonomous multi-model research engine coordinating AI-assisted discovery, synthesis, and verification. Devanik Debnath retains responsibility for the final manuscript, scientific claims, citations, and released artifacts.
\bibliographystyle{plainnat}
\bibliography{paper}
\end{document}
